\section{What a Head Buys, and Where the Gain Comes From}\label{sec:heads}
\Cref{sec:readout} established that the readout reorders the encoders. This section asks what
the readout is actually supplying, using the PCA control to separate two hypotheses that the
supervised rows cannot distinguish on their own.

\begin{table*}[t]
\centering
\caption{The lead encoder (C-RADIOv4-H at $224\times224$) under all four \texttt{market1501/train}-fitted readouts, mAP. Only
the Market-1501 column is in-domain; every other column applies the same Market-fitted head to
a domain whose labels it has never seen. The PCA row is the matched control --- same images,
same output dimension, no labels.}
\label{tab:heads}
\scriptsize
\setlength{\tabcolsep}{3pt}
% generated by tools/gen_tables.py from results/runs — edit the runs, not this
\begin{tabular}{lrrrrrrr}
\toprule
head & Market & MSMT17 & CUHK03 & Occ-REID & CCVID & VRIC & VRAI \\
\midrule
frozen & 0.061 & 0.023 & 0.016 & 0.456 & 0.523 & 0.119 & 0.174 \\
PCA & 0.069 & 0.027 & 0.021 & 0.456 & 0.573 & 0.106 & 0.163 \\
linear & 0.400 & 0.093 & 0.130 & 0.624 & \textbf{0.624} & 0.130 & 0.171 \\
ArcFace & \textbf{0.709} & \textbf{0.165} & \textbf{0.153} & \textbf{0.646} & 0.609 & \textbf{0.191} & \textbf{0.228} \\
\bottomrule
\end{tabular}

\end{table*}

\subsection{The gain is the labels, not the bottleneck}
Every fitted head in this study produces a 512-dimensional embedding from a
higher-dimensional frozen one, so any improvement it shows is confounded between two causes:
the supervision it received, and the dimensionality reduction itself. Without a control, ``the
ArcFace head beats the frozen encoder'' is indistinguishable from ``512 dimensions are enough''.

The PCA row is that control. It is fitted on the same \nHeadImages{} images, produces the same
512-dimensional embedding, and uses none of the \nHeadIdentities{} identity labels. On
the lead encoder it moves Market-1501 from \leadMarketFrozen{} to \leadMarketPCA{}, moves
Occluded-REID from \leadOccludedFrozen{} to \leadOccludedPCA{} --- a change well inside that
dataset's seed noise --- and \emph{loses} on VRAI, \leadVraiFrozen{} to \leadVraiPCA{}. The
bottleneck is worth approximately nothing. The supervised heads on the same features reach
\leadMarketLinear{} and \leadMarketArc{}, so the entire effect reported in \cref{sec:readout}
is attributable to the identity supervision.

\begin{table}[t]
\centering
\caption{The same recipe fitted on each of two training splits, read on both of them, on the
lead encoder. The diagonal is in-domain. Every off-diagonal cell is the same recipe, the same
encoder, the same fitting budget and the same test set, and differs only in where the head was
fitted --- which is what makes this a control rather than a comparison.}
\label{tab:fitsplit}
\small
% generated by tools/gen_tables.py from results/runs — edit the runs, not this
\begin{tabular}{llrr}
\toprule
readout & fitted on & Market & MSMT17 \\
\midrule
frozen & --- & 0.061 & 0.023 \\
\midrule
PCA & Market-1501 & \textbf{0.069} & 0.027 \\
PCA & MSMT17 & 0.063 & \textbf{0.026} \\
\midrule
linear & Market-1501 & \textbf{0.400} & 0.093 \\
linear & MSMT17 & 0.248 & \textbf{0.260} \\
\midrule
ArcFace & Market-1501 & \textbf{0.709} & 0.165 \\
ArcFace & MSMT17 & 0.378 & \textbf{0.488} \\
\bottomrule
\end{tabular}

\end{table}

\subsection{Half of what a head buys is where it was fitted}
\Cref{sec:setup} fixes every head in the preceding tables on \texttt{market1501/train}, which
makes one column of each of them in-domain and every other column a transfer number. That
convention is the usual one, and it is also a confound: it cannot separate \emph{a head helps}
from \emph{a head fitted here helps}, because it never varies the second clause.

\Cref{tab:fitsplit} varies it. The same three recipes, the same encoders, the same budget and
the same seed are refitted on \texttt{msmt17/train} --- \nHeadImagesMsmt{} images of
\nHeadIdentitiesMsmt{} identities --- and each is read on both training domains. Nothing else
changes, and no backbone is re-run: the frozen features of both training splits were already
cached, so the control costs a head fit and not a single forward pass.

The result is a mirror. Averaged over every encoder, ArcFace reads \fitArcfaceMarketHere{} mAP
on Market-1501 when it was fitted there and \fitArcfaceMarketElse{} when the identical recipe
was fitted on MSMT17; on MSMT17 the same two heads read \fitArcfaceMsmtHere{} and
\fitArcfaceMsmtElse{}. The swing is \fitSwingArcfaceMarket{} in one direction and
\fitSwingArcfaceMsmt{} in the other. A gap visible in one direction only could be a property of
MSMT17; one that appears in both directions at the same size is a property of fitting.

It is monotone in supervision, and the PCA control is again what makes that legible. The same
swap moves PCA by \fitSwingPcaMarket{}, the linear head by \fitSwingLinearMarket{} and ArcFace
by \fitSwingArcfaceMarket{}. An unsupervised rotation of the same features barely encodes which
domain it was fitted on; an angular-margin classifier over identity labels encodes it most.

This sharpens \cref{sec:readout}'s claim rather than weakening it. A head that has never seen a
Market-1501 label still lifts Market-1501 mAP from \fitFrozenMarket{} to
\fitArcfaceMarketElse{}, so most of what the readout supplies is a transferable metric geometry
and not memorised identities; fitting it in-domain then roughly doubles that again. Of the
total gain over the frozen encoder, \fitShareArcfaceMarket{} survives being fitted elsewhere on
Market-1501 and \fitShareArcfaceMsmt{} does on MSMT17. The readout still determines the
ranking, as \cref{sec:readout} reports. What the fit domain determines is how much of the
readout a deployment somewhere else gets to keep.

One asymmetry in \cref{tab:fitsplit} carries less than it appears to. \texttt{msmt17/train} is
\nHeadImagesMsmt{} images against Market-1501's \nHeadImages{}, so which of the two splits
yields the better \emph{transfer} head is confounded with how much data each head saw. The
mirror above is read within a column and does not rest on that comparison; a matched-budget
version of it is left to future work.

\subsection{A frozen backbone plus one affine layer is a real system}
The best in-domain configuration in this study reaches \bestMarketMAP{} mAP and
\bestMarketRone{} Rank-1 on \texttt{market1501/official@1}. This is the first number in the
study that may be placed beside a published one, because it is produced the way published
numbers are: fitted on the official training split, evaluated on the official protocol, with
the test identities disjoint from the training identities. It does not beat a tuned
specialist. It costs one forward pass over \nHeadImages{} cached crops plus well under a
minute of head fitting on a single consumer GPU, with no backbone gradient at any point.

\subsection{Angular margin, and where it fails}
ArcFace beats a plain linear head almost everywhere in this study, and by the largest margin
where the head was fitted --- \leadMarketLinear{} to \leadMarketArc{} on Market-1501. This is
expected: an angular margin optimises the cosine geometry the retrieval metric reads, and a
softmax cross-entropy head does not. The exceptions are informative and are concentrated in one
encoder family; \cref{tab:margin} shows them, and \cref{sec:ablation} interprets them.

\begin{table*}[t]
\centering
\caption{Relative change in mAP from replacing the linear head with the ArcFace head, same
features and same budget. Positive is the expected direction. The losses cluster in DINOv3's
feature space on transfer sets, and both DINOv3 probes reach a perfect training accuracy, so
they are not fit failures.}
\label{tab:margin}
\scriptsize
\setlength{\tabcolsep}{3pt}
% generated by tools/gen_tables.py from results/runs — edit the runs, not this
\begin{tabular}{lrrrrrrr}
\toprule
encoder & Market & MSMT17 & CUHK03 & Occ-REID & CCVID & VRIC & VRAI \\
\midrule
C-RADIOv4-H & +77\% & +76\% & +17\% & +3\% & \loss{-2\%} & +47\% & +33\% \\
C-RADIOv4-SO400M & +89\% & +77\% & +36\% & +7\% & +3\% & +41\% & +41\% \\
SigLIP2-g & +51\% & +28\% & +5\% & +3\% & +1\% & +39\% & +38\% \\
SigLIP2-SO400M & +78\% & +32\% & +16\% & +5\% & +2\% & +56\% & +36\% \\
DINOv3-H+ & +105\% & +70\% & \loss{-20\%} & \loss{-3\%} & \loss{-4\%} & +10\% & +32\% \\
DINOv3-L & +101\% & +53\% & +1\% & \loss{-9\%} & +4\% & +2\% & +17\% \\
CLIP-B/16 & +103\% & +4\% & \loss{-19\%} & +0\% & +14\% & +25\% & +11\% \\
\bottomrule
\end{tabular}

\end{table*}

\subsection{The head transfers out of its object class}
The same head fitted on \nHeadIdentities{} pedestrian identities improves aerial \emph{vehicle}
retrieval, from \leadVraiFrozen{} frozen to \leadVraiArc{} with ArcFace, having never seen a
vehicle label. Since the matched PCA control \emph{falls} on that dataset, this is not a
dimensionality effect either. What the head learned is not ``what a person looks like'' but a
metric geometry that instance discrimination reuses across object classes --- which is the
most useful practical finding in this section and the one that most needs a larger transfer
grid before it is leaned on.
